\section{Conclusion and future work}
\label{sec:conclusion}

Two research communities study one failure. An LLM agent reads content it
should not trust, and it cannot tell which content that is. The security
community calls the content an attack and builds defenses. The machine learning
community calls it a mechanism and builds fixes. We placed 984 papers from both
sides into one taxonomy of channel, intent and consequence, and found that the
separation is measurable: about one paper in ten cites the other side, 14 of
529 defenses were validated against both threat models, and 99.1\% of tested
defense-mechanism pairs stay inside a single track.

We then used the taxonomy to find pairs that should have been run and were not,
and ran 81 of them on the authors' released artifacts. The central finding is
that a popular family of fixes, those that repair accidental context failure by
shifting trust from the model's parameters toward the retrieved text, becomes a
liability under deliberate attack, for the reason that is written into its own
design. The finding has a measured scope condition and a counterexample, so it
is a design error rather than a law.

A second result cuts across the first. Protection depends on the \emph{form}
of the planted passage, not only on whether an attacker is present. Running the
same defenses against an optimized adversarial passage and against a fluent
sentence that asserts a plausible falsehood gives different answers, and the
one defense that resists the first falls below a no-retrieval baseline against
the second. The fluent payload needs no gradient access and no surrogate
retriever, and it damages an undefended model roughly twice as hard. An
evaluation built only around optimized payloads overstates how well a system
holds up.

\paragraph{Future work.}
Four directions follow directly.

\textbf{Close the supply-chain and cross-modal gaps.} These channels have fewer
than one named defense per named mechanism and are the routes by which agents
accept third-party capability. They are the clearest place to build rather than
to measure.

\textbf{Separate the intervention-point finding from its confound.} Our case
studies suggest that whether a defense generalizes tracks where in the pipeline
it intervenes. The ingestion defenses that generalize are also the ones whose
harnesses model a text-classification victim, while the execution defenses we
could not evaluate are the agent-shaped ones. Running all three intervention
points inside one agent harness would resolve this, and that work is in
progress.

\textbf{Separate recall from copying.} Under false evidence, the planted
passage asserts the same string the model would have recalled, so our
parametric-answer metric cannot distinguish a model that remembered from a
model that copied. An arm that plants the gold answer instead would separate
them and would sharpen the mechanism claim in Section~\ref{sec:results}.

\textbf{Build the missing evaluation.} The deepest problem this paper documents
is not any single defense. It is that the field's reported outcomes are
conditioned on who ran the experiment, and that nothing in current practice
corrects for it. A shared benchmark that pairs defenses with attacks selected
by a taxonomy rather than by either author would make transfer a measured
property rather than an assumed one.
